\ifdefined\gkver\else\def\gkver{student}\fi
\documentclass[\gkver]{gaokaozhenti}
\begin{document}

\chapter{2005年江苏理综}
%% paperType: 理综物理部分

\begin{enumerate}
\item
%% number: 1
%% paperName: 2005年江苏理综第 1 题 4 分
%% typeId: 2
%% type: 多选题
%% chapter: 原子和原子核
%% point: 人工核反应
%% method: 无
%% score: 4
%% degree: 650
%% duplicateId: 0
%% body:
下列核反应或核衰变方程中，符号“X”表示中子的是\xzanswer{AC}

\fourchoices[answer=AC]
{\ce{^9_4Be + ^4_2He -> ^12_6C + X}}
{\ce{^14_7N + ^4_2He -> ^17_8O + X}}
{\ce{^204_80Hg + ^1_0n -> ^202_78Pt + 2^1_1H + X}}
{\ce{^239_92U -> ^239_93Np + X}}

%% answer: AC
%% memo:
\memoanswer{
本题原卷及材料中未提供详解，待补充。
}

\item
%% number: 2
%% paperName: 2005年江苏理综第 2 题 4 分
%% typeId: 2
%% type: 多选题
%% chapter: 光学
%% point: 光电效应
%% method: 无
%% score: 4
%% degree: 650
%% duplicateId: 0
%% body:
为了强调物理学对当今社会的重要作用并纪念爱因斯坦，2004 年联合国第 58 次大会把 2005 年定为国际物理年。爱因斯坦在 100 年前发表了 5 篇重要论文，内容涉及狭义相对论、量子论和统计物理学，对现代物理学的发展作出了巨大贡献。某人学了有关的知识后，有如下理解，其中正确的是\xzanswer{BD}

\fourchoices[answer=BD]
{所谓布朗运动就是液体分子的无规则运动}
{光既具有波动性，又具有粒子性}
{在光电效应的实验中，入射光强度增大，光电子的最大初动能随之增大}
{质能方程表明：物体具有的能量与它的质量有简单的正比关系}

%% answer: BD
%% memo:
\memoanswer{
本题原卷及材料中未提供详解，待补充。
}

\item
%% number: 3
%% paperName: 2005年江苏理综第 3 题 4 分
%% typeId: 2
%% type: 多选题
%% chapter: 电场
%% point: 等势面
%% method: 无
%% score: 4
%% degree: 650
%% duplicateId: 0
%% body:
根据 $\alpha$ 粒子散射实验，卢瑟福提出了原子的核式结构模型。图中虚线表示原子核所形成的电场的等势线，实线表示一个 $\alpha$ 粒子的运动轨迹。在 $\alpha$ 粒子从 a 运动到 b、再运动到 c 的过程中，下列说法中正确的是\xzanswer{C}

\onepicture[w=3.6cm]{figs/03.png}

\fourchoices[answer=C]
{动能先增大，后减小}
{电势能先减小，后增大}
{电场力先做负功，后做正功，总功等于零}
{加速度先变小，后变大}

%% answer: C
%% memo:
\memoanswer{
本题原卷及材料中未提供详解，待补充。
}

\item
%% number: 4
%% paperName: 2005年江苏理综第 4 题 4 分
%% typeId: 2
%% type: 多选题
%% chapter: 气体性质
%% point: 阿伏伽德罗常数
%% method: 无
%% score: 4
%% degree: 650
%% duplicateId: 0
%% body:
某气体的摩尔质量为 $M$，摩尔体积为 $V$，密度为 $\rho$，每个分子的质量和体积分别为 $m$ 和 $V_{0}$，则阿伏加德罗常数 $N_{A}$ 可表示为\xzanswer{BC}

\fourchoices[answer=BC]
{$N_{A}=\frac{V}{V_{0}}$}
{$N_{A}=\frac{\rho V}{m}$}
{$N_{A}=\frac{M}{m}$}
{$N_{A}=\frac{M}{\rho V_{0}}$}

%% answer: BC
%% memo:
\memoanswer{
本题原卷及材料中未提供详解，待补充。
}

\item
%% number: 5
%% paperName: 2005年江苏理综第 5 题 4 分
%% typeId: 1
%% type: 单选题
%% chapter: 曲线运动
%% point: 人造地球卫星
%% method: 无
%% score: 4
%% degree: 650
%% duplicateId: 0
%% body:
某人造卫星运动的轨道可近似看作是以地心为中心的圆。由于阻力作用，人造卫星到地心的距离从 $r_{1}$ 慢慢变到 $r_{2}$，用 $E_{k1}$、$E_{k2}$ 分别表示卫星在这两个轨道上的动能，则\xzanswer{B}

\fourchoices[answer=B]
{$r_{1}<r_{2}$，$E_{k1}<E_{k2}$}
{$r_{1}>r_{2}$，$E_{k1}<E_{k2}$}
{$r_{1}<r_{2}$，$E_{k1}>E_{k2}$}
{$r_{1}>r_{2}$，$E_{k1}>E_{k2}$}

%% answer: B
%% memo:
\memoanswer{
本题原卷及材料中未提供详解，待补充。
}

\item
%% number: 6
%% paperName: 2005年江苏理综第 6 题 4 分
%% typeId: 1
%% type: 单选题
%% chapter: 光学
%% point: 光子说
%% method: 无
%% score: 4
%% degree: 650
%% duplicateId: 0
%% body:
在中子衍射技术中，常利用热中子研究晶体的结构，因为热中子的德布罗意波长与晶体中原子间距相近。已知中子质量 $m=1.67\times10^{-27}\Ukg$，普朗克常量 $h=6.63\times10^{-34}\UJcdots$，可以估算德布罗意波长 $\lambda=1.82\times10^{-10}\Um$ 的热中子动能的数量级为\xzanswer{C}

\fourchoices[answer=C]
{$10^{-17}\UJ$}
{$10^{-19}\UJ$}
{$10^{-21}\UJ$}
{$10^{-24}\UJ$}

%% answer: C
%% memo:
\memoanswer{
本题原卷及材料中未提供详解，待补充。
}

\item
%% number: 7
%% paperName: 2005年江苏理综第 7 题 4 分
%% typeId: 1
%% type: 单选题
%% chapter: 气体性质
%% point: 分子力
%% method: 无
%% score: 4
%% degree: 650
%% duplicateId: 0
%% body:
下列关于分子力和分子势能的说法中，正确的是\xzanswer{C}

\fourchoices[answer=C]
{当分子力表现为引力时，分子力和分子势能总是随分子间距离的增大而增大}
{当分子力表现为引力时，分子力和分子势能总是随分子间距离的增大而减小}
{当分子力表现为斥力时，分子力和分子势能总是随分子间距离的减小而增大}
{当分子力表现为斥力时，分子力和分子势能总是随分子间距离的减小而减小}

%% answer: C
%% memo:
\memoanswer{
本题原卷及材料中未提供详解，待补充。
}

\item
%% number: 8
%% paperName: 2005年江苏理综第 8 题 4 分
%% typeId: 2
%% type: 多选题
%% chapter: 机械波
%% point: 波的图象
%% method: 无
%% score: 4
%% degree: 650
%% duplicateId: 0
%% body:
一列简谐横波沿 $x$ 轴传播。$T=0$ 时的波形如图所示，质点 A 与质点 B 相距 $1\Um$，A 点速度沿 $y$ 轴正方向；$t=0.02\Us$ 时，质点 A 第一次到达正向最大位移处。由此可知\xzanswer{AB}

\onepicture[w=6.0cm]{figs/08.png}

\fourchoices[answer=AB]
{此波的传播速度为 $25\Ums$}
{此波沿 $x$ 轴负方向传播}
{从 $t=0$ 时起，经过 $0.04\Us$，质点 A 沿波传播方向迁移了 $1\Um$}
{在 $t=0.04\Us$ 时，质点 B 处在平衡位置，速度沿 $y$ 轴负方向}

%% answer: AB
%% memo:
\memoanswer{
本题原卷及材料中未提供详解，待补充。
}

\item
%% number: 9
%% paperName: 2005年江苏理综第 9 题 4 分
%% typeId: 1
%% type: 单选题
%% chapter: 气体性质
%% point: 热力学第一定律
%% method: 无
%% score: 4
%% degree: 450
%% duplicateId: 0
%% body:
分别以 $p$、$V$、$T$ 表示气体的压强、体积、温度。一定质量的理想气体，其初始状态表示为 $(p_{0}, V_{0}, T_{0})$。若分别经历如下两种变化过程：

①从 $(p_{0}, V_{0}, T_{0})$ 变为 $(p_{1}, V_{1}, T_{1})$ 的过程中，温度保持不变 $(T_{1}=T_{0})$；

②从 $(p_{0}, V_{0}, T_{0})$ 变为 $(p_{2}, V_{2}, T_{2})$ 的过程中，既不吸热，也不放热。

在上述两种变化过程中，如果 $V_{1}=V_{2}>V_{0}$，则\xzanswer{A}

\fourchoices[answer=A]
{$p_{1}>p_{2}$，$T_{1}>T_{2}$}
{$p_{1}>p_{2}$，$T_{1}<T_{2}$}
{$p_{1}<p_{2}$，$T_{1}<T_{2}$}
{$p_{1}<p_{2}$，$T_{1}>T_{2}$}

%% answer: A
%% memo:
\memoanswer{
本题原卷及材料中未提供详解，待补充。
}

\item
%% number: 10
%% paperName: 2005年江苏理综第 10 题 4 分
%% typeId: 1
%% type: 单选题
%% chapter: 功和能
%% point: 功
%% method: 无
%% score: 4
%% degree: 450
%% duplicateId: 0
%% body:
如图所示，固定的光滑竖直杆上套着一个滑块，用轻绳系着滑块绕过光滑的定滑轮，以大小恒定的拉力 $F$ 拉绳，使滑块从 A 点起由静止开始上升。若从 A 点上升至 B 点和从 B 点上升至 C 点的过程中拉力 $F$ 做的功分别为 $W_{1}$、$W_{2}$，滑块经 B、C 两点时的动能分别为 $E_{kB}$、$E_{kC}$，图中 AB$=$BC，则一定有\xzanswer{A}

\onepicture[w=5.0cm]{figs/10.png}

\fourchoices[answer=A]
{$W_{1}>W_{2}$}
{$W_{1}<W_{2}$}
{$E_{kB}>E_{kC}$}
{$E_{kB}<E_{kC}$}

%% answer: A
%% memo:
\memoanswer{
本题原卷及材料中未提供详解，待补充。
}

\item
%% number: 11
%% paperName: 2005年江苏理综第 11 题 10 分
%% typeId: 4
%% type: 实验题
%% chapter: 力物体的平衡
%% point: 胡克定律
%% method: 数据归纳
%% score: 10
%% degree: 650
%% duplicateId: 0
%% body:
某同学用如图所示装置做探究弹力和弹簧伸长关系的实验。他先测出不挂砝码时弹簧下端指针所指的标尺刻度，然后在弹簧下端挂上砝码，并逐个增加砝码，测出指针所指的标尺刻度，所得数据列表如下：（重力加速度 $g=9.8\Umsq$）

\begin{center}
\begin{tabular}{|c|c|c|c|c|c|c|c|c|}
\hline
砝码质量 $m/\times10^{2}\Ug$ & 0 & 1.00 & 2.00 & 3.00 & 4.00 & 5.00 & 6.00 & 7.00 \\
\hline
标尺刻度 $x/10^{-2}\Um$ & 15.00 & 18.94 & 22.82 & 26.78 & 30.66 & 34.60 & 42.00 & 54.50 \\
\hline
\end{tabular}
\end{center}

\begin{enumerate}
	\item
	根据所测数据，在答题卡的坐标纸上作出弹簧指针所指的标尺刻度与砝码质量的关系曲线。

	\item
	根据所测得的数据和关系曲线可以判断，在 \tkanswer[2cm]{$0\sim4.9$} 范围内弹力大小与弹簧伸长关系满足胡克定律。这种规格弹簧的劲度系数为 \tkanswer[2cm]{25.0} $\UNm$。
\end{enumerate}

\twopicture[showlabel=false, h=4.6cm, align=right]
{figs/11a.png}
{figs/11b.png}

%% answer: 
\jdanswer{
\begin{enumerate}
	\item
	\onepicture[w=8cm]{figs/11_an.png}

	\item
	$0\sim4.9$，25.0
\end{enumerate}
}
%% memo:
\memoanswer{
本题原卷及材料中未提供详解，待补充。
}

\item
%% number: 12
%% paperName: 2005年江苏理综第 12 题 12 分
%% typeId: 4
%% type: 实验题
%% chapter: 电路
%% point: 测电动势与内阻
%% method: 无
%% score: 12
%% degree: 450
%% duplicateId: 0
%% body:
将满偏电流 $I_{g}=300\UuA$、内阻未知的电流表 G 改装成电压表并进行核对。

\begin{enumerate}
	\item
	利用如图所示的电路测量电流表 G 的内阻（图中电源的电动势 $E=4\UV$）：先闭合 $S_{1}$，调节 $R$，使电流表指针偏转到满刻度；再闭合 $S_{2}$，保持 $R$ 不变，调节 $R'$，使电流表指针偏转到满刻度的 2/3，读出此时 $R'$ 的阻值为 $200\UO$，则电流表内阻的测量值 $R_{g}=$\tkanswer[2cm]{100} $\UO$。

	\item
	将该表改装成量程为 $3\UV$ 的电压表，需 \tkanswer[1.5cm]{串联}（填“串联”或“并联”）阻值为 $R_{0}=$\tkanswer[2cm]{9900} $\UO$ 的电阻。

	\item
	把改装好的电压表与标准电压表进行核对，试在答题卡上画出实验电路图和实物连接图。
\end{enumerate}

\onepicture[w=5.2cm, align=right]{figs/12a.png}

\onepicture[w=10.5cm, align=right]{figs/12b.png}

%% answer: 
\jdanswer{
\begin{enumerate}
	\item
	100

	\item
	串联，9900

	\item
	\onepicture[w=5.5cm]{figs/12_an1.png}

	\onepicture[w=8.5cm]{figs/12_an2.png}
\end{enumerate}
}
%% memo:
\memoanswer{
本题原卷及材料中未提供详解，待补充。
}

\item
%% number: 13
%% paperName: 2005年江苏理综第 13 题 14 分
%% typeId: 6
%% type: 计算题
%% chapter: 曲线运动
%% point: 平抛运动
%% method: 无
%% score: 14
%% degree: 650
%% duplicateId: 0
%% body:
A、B 两小球同时从距地面高为 $h=15\Um$ 处的同一点抛出，初速度大小均为 $v_{0}=10\Ums$。A 球竖直向下抛出，B 球水平抛出，空气阻力不计。求：

\begin{enumerate}
	\item
	A 球经多长时间落地？

	\item
	A 球落地时，A、B 两球间的距离是多少？
\end{enumerate}

%% answer: 
\jdanswer{
\begin{enumerate}
	\item
	$t=1\Us$

	\item
	$L=10\sqrt{2}\Um$
\end{enumerate}
}
%% memo:
\memoanswer{
（1）A 球做竖直下抛运动：$h=v_{0}t+\frac{1}{2}gt^{2}$

将 $h=15\Um$、$v_{0}=10\Ums$ 代入，可得：$t=1\Us$。

（2）B 球做平抛运动：

$x=v_{0}t$，$y=\frac{1}{2}gt^{2}$

将 $v_{0}=10\Ums$、$t=1\Us$ 代入，可得：$x=10\Um$，$y=5\Um$。

此时 A 球与 B 球的距离 $L$ 为：$L=\sqrt{x^{2}+(h-y)^{2}}$

将 $x$、$y$、$h$ 代入，得：$L=10\sqrt{2}\Um$。
}

\item
%% number: 14
%% paperName: 2005年江苏理综第 14 题 12 分
%% typeId: 6
%% type: 计算题
%% chapter: 电路
%% point: 输出功率极值
%% method: 极值
%% score: 12
%% degree: 650
%% duplicateId: 0
%% body:
如图所示，$R$ 为电阻箱，V 为理想电压表。当电阻箱读数为 $R_{1}=2\UO$ 时，电压表读数为 $U_{1}=4\UV$；当电阻箱读数为 $R_{2}=5\UO$ 时，电压表读数为 $U_{2}=5\UV$。求：

\begin{enumerate}
	\item
	电源的电动势 $E$ 和内阻 $r$；

	\item
	当电阻箱 $R$ 读数为多少时，电源的输出功率最大？最大值 $P_{m}$ 为多少？
\end{enumerate}

\onepicture[w=5.0cm, align=right]{figs/14.png}

%% answer: 
\jdanswer{
\begin{enumerate}
	\item
	$E=6\UV$，$r=1\UO$

	\item
	$R=r=1\UO$ 时 $P$ 有最大值，$P_{m}=\frac{E^{2}}{4r}=9\UW$
\end{enumerate}
}
%% memo:
\memoanswer{
（1）由闭合电路欧姆定律：$E=U_{1}+\frac{U_{1}}{R_{1}}r$，$E=U_{2}+\frac{U_{2}}{R_{2}}r$

联立上式并代入数据解得：$E=6\UV$，$r=1\UO$。

（2）由电功率表达式：$P=\frac{E^{2}R}{(R+r)^{2}}$

将上式变形为：$P=\frac{E^{2}}{\frac{(R-r)^{2}}{R}+4r}$

由上式可知 $R=r=1\UO$ 时 $P$ 有最大值，$P_{m}=\frac{E^{2}}{4r}=9\UW$。
}

\item
%% number: 15
%% paperName: 2005年江苏理综第 15 题 14 分
%% typeId: 6
%% type: 计算题
%% chapter: 光学
%% point: 光的干涉
%% method: 无
%% score: 14
%% degree: 650
%% duplicateId: 0
%% body:
1801 年，托马斯$\cdot$杨用双缝干涉实验研究了光波的性质。1834 年，洛埃利用单面镜同样得到了杨氏干涉的结果（称洛埃镜实验）。

\begin{enumerate}
	\item
	洛埃镜实验的基本装置如图所示，$S$ 为单色光源，$M$ 为一平面镜。试用平面镜成像作图法在答题卡上画出 $S$ 经平面镜反射后的光与直接发出的光在光屏上相交的区域。

	\item
	设光源 $S$ 到平面镜的垂直距离和到光屏的垂直距离分别为 $a$ 和 $L$，光的波长为 $\lambda$，在光屏上形成干涉条纹。写出相邻两条亮纹（或暗纹）间距离 $\Delta x$ 的表达式。
\end{enumerate}

\onepicture[w=7.6cm, align=right]{figs/15.png}

%% answer: 
\jdanswer{
\begin{enumerate}
	\item
	\onepicture[w=7.6cm]{figs/15_an.png}

	\item
	$\Delta x=\frac{L\lambda}{2a}$
\end{enumerate}
}
%% memo:
\memoanswer{
（1）光路图如图所示。

（2）双缝干涉相邻两条亮纹（或暗纹）间的距离 $\Delta x=\frac{L\lambda}{d}$，洛埃镜的等效双缝间距 $d=2a$，所以

$\Delta x=\frac{L\lambda}{2a}$。
}

\item
%% number: 16
%% paperName: 2005年江苏理综第 16 题 16 分
%% typeId: 6
%% type: 计算题
%% chapter: 电磁感应
%% point: 电磁感应能量分析
%% method: 无
%% score: 16
%% degree: 450
%% duplicateId: 0
%% body:
如图所示，固定的水平光滑金属导轨，间距为 $L$，左端接有阻值为 $R$ 的电阻，处在方向竖直、磁感应强度为 $B$ 的匀强磁场中，质量为 $m$ 的导体棒与固定弹簧相连，放在导轨上，导轨与导体棒的电阻均可忽略。初始时刻，弹簧恰处于自然长度，导体棒具有水平向右的初速度 $v_{0}$，在沿导轨往复运动的过程中，导体棒始终与导轨垂直并保持良好接触。

\begin{enumerate}
	\item
	求初始时刻导体棒受到的安培力。

	\item
	若导体棒从初始时刻到速度第一次为零时，弹簧的弹性势能为 $E_{p}$，则这一过程中安培力所做的功 $W_{1}$ 和电阻 $R$ 上产生的焦耳热 $Q_{1}$ 分别为多少？

	\item
	导体棒往复运动，最终将静止于何处？从导体棒开始运动直到最终静止的过程中，电阻 $R$ 上产生的焦耳热 $Q$ 为多少？
\end{enumerate}

\onepicture[w=6.2cm, align=right]{figs/16.png}

%% answer: 
\jdanswer{
\begin{enumerate}
	\item
	$F=\frac{B^{2}L^{2}v_{0}}{R}$，方向水平向左

	\item
	$W_{1}=E_{p}-\frac{1}{2}mv_{0}^{2}$，$Q_{1}=\frac{1}{2}mv_{0}^{2}-E_{p}$

	\item
	棒最终静止于初始位置，$Q=\frac{1}{2}mv_{0}^{2}$
\end{enumerate}
}
%% memo:
\memoanswer{
（1）初始时刻棒中感应电动势 $E=BLv_{0}$

棒中感应电流 $I=\frac{E}{R}$

作用于棒上的安培力 $F=ILB$

联立得 $F=\frac{B^{2}L^{2}v_{0}}{R}$，安培力方向水平向左。

（2）由功和能的关系，得安培力做功 $W_{1}=E_{p}-\frac{1}{2}mv_{0}^{2}$，电阻 $R$ 上产生的焦耳热 $Q_{1}=\frac{1}{2}mv_{0}^{2}-E_{p}$。

（3）由能量转化及平衡条件等，可判断：棒最终静止于初始位置，$Q=\frac{1}{2}mv_{0}^{2}$。
}

\item
%% number: 17
%% paperName: 2005年江苏理综第 17 题 16 分
%% typeId: 6
%% type: 计算题
%% chapter: 磁场
%% point: 洛伦兹力
%% method: 无
%% score: 16
%% degree: 450
%% duplicateId: 0
%% body:
如图所示，M、N 为两块带等量异种电荷的平行金属板，$S_{1}$、$S_{2}$ 为板上正对的小孔，N 板右侧有两个宽度均为 $d$ 的匀强磁场区域，磁感应强度大小均为 $B$，方向分别垂直于纸面向外和向里，磁场区域右侧有一个荧光屏，取屏上与 $S_{1}$、$S_{2}$ 共线的 O 点为原点，向上为正方向建立 $x$ 轴。M 板左侧电子枪发射出的热电子经小孔 $S_{1}$ 进入两板间，电子的质量为 $m$，电荷量为 $e$，初速度可以忽略。

\begin{enumerate}
	\item
	当两板间电势差为 $U_{0}$ 时，求从小孔 $S_{2}$ 射出的电子的速度 $v_{0}$。

	\item
	求两金属板间电势差 $U$ 在什么范围内，电子不能穿过磁场区域而打到荧光屏上。

	\item
	若电子能够穿过磁场区域而打到荧光屏上，试在答题卡的图上定性地画出电子运动的轨迹。

	\item
	求电子打到荧光屏上的位置坐标 $x$ 和金属板间电势差 $U$ 的函数关系。
\end{enumerate}

\onepicture[w=8.0cm, align=right]{figs/17.png}

%% answer: 
\jdanswer{
\begin{enumerate}
	\item
	$v_{0}=\sqrt{\frac{2eU_{0}}{m}}$

	\item
	$U<\frac{d^{2}eB^{2}}{2m}$

	\item
	\onepicture[w=8.0cm]{figs/17_an.png}

	\item
	$x=\frac{2(\sqrt{2emU}-\sqrt{2emU-d^{2}e^{2}B^{2}})}{eB}$（$U\geqslant\frac{d^{2}eB^{2}}{2m}$）
\end{enumerate}
}
%% memo:
\memoanswer{
（1）根据动能定理，得 $eU_{0}=\frac{1}{2}mv_{0}^{2}$，由此可解得 $v_{0}=\sqrt{\frac{2eU_{0}}{m}}$。

（2）欲使电子不能穿过磁场区域而打到荧光屏上，应有 $r=\frac{mv}{eB}<d$

而 $eU=\frac{1}{2}mv^{2}$，由此即可解得 $U<\frac{d^{2}eB^{2}}{2m}$。

（3）电子穿过磁场区域而打到荧光屏上时运动的轨迹如图所示。

（4）若电子在磁场区域做圆周运动的轨道半径为 $r$，穿过磁场区域打到荧光屏上的位置坐标为 $x$，则由（3）中的轨迹图可得 $x=2r-2\sqrt{r^{2}-d^{2}}$。注意到 $r=\frac{mv}{eB}$ 和 $eU=\frac{1}{2}mv^{2}$，所以电子打到荧光屏上的位置坐标 $x$ 和金属板间电势差 $U$ 的函数关系为

$x=\frac{2(\sqrt{2emU}-\sqrt{2emU-d^{2}e^{2}B^{2}})}{eB}$（$U\geqslant\frac{d^{2}eB^{2}}{2m}$）。
}

\item
%% number: 18
%% paperName: 2005年江苏理综第 18 题 16 分
%% typeId: 6
%% type: 计算题
%% chapter: 运动和力
%% point: 动量守恒定律
%% method: 无
%% score: 16
%% degree: 250
%% duplicateId: 0
%% body:
如图所示，三个质量均为 $m$ 的弹性小球用两根长均为 $L$ 的轻绳连成一条直线而静止在光滑水平面上。现给中间的小球 B 一个水平初速度 $v_{0}$，方向与绳垂直。小球相互碰撞时无机械能损失，轻绳不可伸长。求：

\begin{enumerate}
	\item
	当小球 A、C 第一次相碰时，小球 B 的速度。

	\item
	当三个小球再次处在同一直线上时，小球 B 的速度。

	\item
	运动过程中小球 A 的最大动能 $E_{kA}$ 和此时两根绳的夹角 $\theta$。

	\item
	当三个小球处在同一直线上时，绳中的拉力 $F$ 的大小。
\end{enumerate}

\onepicture[w=7.0cm, align=right]{\includesvg{figs/18}}

%% answer: 
\jdanswer{
\begin{enumerate}
	\item
	$v_{B}=\frac{1}{3}v_{0}$

	\item
	$v_{B}=-\frac{1}{3}v_{0}$

	\item
	$E_{kA}=\frac{1}{4}mv_{0}^{2}$，$\theta=90^{\circ}$

	\item
	$F=\frac{mv_{0}^{2}}{L}$
\end{enumerate}
}
%% memo:
\memoanswer{
（1）设小球 A、C 第一次相碰时，小球 B 的速度为 $v_{B}$，考虑到对称性及绳的不可伸长特性，小球 A、C 沿小球 B 初速度方向的速度也为 $v_{B}$，由动量守恒定律，得

$mv_{0}=3mv_{B}$

由此解得 $v_{B}=\frac{1}{3}v_{0}$。

（2）当三个小球再次处在同一直线上时，则由动量守恒定律和机械能守恒定律，得

$mv_{0}=mv_{B}+2mv_{A}$

$\frac{1}{2}mv_{0}^{2}=\frac{1}{2}mv_{B}^{2}+mv_{A}^{2}$

解得 $v_{B}=-\frac{1}{3}v_{0}$，$v_{A}=\frac{2}{3}v_{0}$（三球再次处于同一直线）

$v_{B}=v_{0}$，$v_{A}=0$（初始状态，舍去）

所以，三个小球再次处在同一直线上时，小球 B 的速度为 $v_{B}=-\frac{1}{3}v_{0}$（负号表明与初速度反向）。

（3）当小球 A 的动能最大时，小球 B 的速度为零。设此时小球 A、C 的速度大小为 $u$，两根绳间的夹角为 $\theta$（如图），则仍由动量守恒定律和机械能守恒定律，得：

$mv_{0}=2mu\sin\frac{\theta}{2}$

$\frac{1}{2}mv_{0}^{2}=mu^{2}$

另外，$E_{kA}=\frac{1}{2}mu^{2}$

\onepicture[w=4.5cm]{figs/18_an.png}

由此可解得，小球 A 的最大动能为 $E_{kA}=\frac{1}{4}mv_{0}^{2}$，此时两根绳间夹角为 $\theta=90^{\circ}$。

（4）小球 A、C 均以半径 $L$ 绕小球 B 做圆周运动，当三个小球处在同一直线上时，以小球 B 为参考系（小球 B 的加速度为 0，为惯性参考系），小球 A（C）相对于小球 B 的速度均为 $v=|v_{A}-v_{B}|=v_{0}$。所以，此时绳中拉力大小为 $F=\frac{mv_{0}^{2}}{L}$。
}

\end{enumerate}

\end{document}
